% =====================================================================
%  ENTRADA DE CARPETAS: mete todos los .tex de una carpeta, en orden.
%
%  \inputcarpeta{ruta}
%
%  Mete los .tex de esa carpeta Y de sus subcarpetas, en orden.
%
%  Por ejemplo, si en la carpeta "Tema 2" tienes:
%     1.1-maquina-y-servidor.tex
%     1.1.1-localhost-y-servidor.tex
%     1.2-formas-de-desplegar.tex
%     1.2.1-tipos-de-despliegue.tex
%     Tema2.1/
%       1.1-maquina-y-servidor.tex
%       1.2-formas-de-desplegar.tex
%     Tema2.2/
%       2.1-apartado.tex
%
%  con \inputcarpeta{Tema 2} los mete TODOS, ordenados (1.1, 1.2,
%  2.1...) y asi no escribes una linea \input por cada archivo.
%
%  Las rutas se cuentan desde la RAIZ del documento (la carpeta donde
%  compilas). Necesita -shell-escape (ya activado en VS Code).
% =====================================================================
% ---------------------------------------------------------------------
%  Etiqueta de la carpeta para el fichero temporal.
%
%  OJO: antes el temporal se llamaba solo /tmp/apuntes-entradas-<jobname>-<n>.tex.
%  Dos documentos distintos con el mismo \jobname (por ejemplo
%  DAW/Tema 2/tema2.tex y DIW/Tema 2/tema2.tex, los dos "tema2")
%  escribian EL MISMO fichero: si se compilaban a la vez (autoguardado +
%  play, dos ventanas...), uno leia la lista de apartados del otro y salia
%  "File `Tema 2.1/...tex' not found".
%  Ahora el nombre incluye la carpeta, asi que cada documento tiene el suyo.
% ---------------------------------------------------------------------
\ExplSyntaxOn
\cs_new_protected:Npn \apuntesetiquetaset #1
  {
    \tl_set:Nn \l_tmpa_tl {#1}
    % todo lo que no sea letra o numero pasa a "-":
    % "Tema 2.1" -> "Tema-2-1", "Tema2.1" -> "Tema2-1"
    \regex_replace_all:nnN { [^A-Za-z0-9]+ } { - } \l_tmpa_tl
  }
\cs_new:Npn \apuntesetiqueta { \tl_use:N \l_tmpa_tl }
\ExplSyntaxOff

\newcount\apuntescontador
\newcommand{\inputcarpeta}[1]{%
  \begingroup
    % \advance dentro de un grupo es local: con \global el contador
    % sube de verdad y cada carpeta usa su propio fichero temporal.
    \global\advance\apuntescontador by 1
    % la carpeta en la que estamos: la usan \imagen y \filaimagen
    % para encontrar las imagenes (imagenes/ de este tema).
    \def\apuntescarpeta{#1}%
    \apuntesetiquetaset{#1}%
    \edef\apuntesentradasfile{/tmp/apuntes-entradas-\jobname-\the\apuntescontador-\apuntesetiqueta.tex}%
    \immediate\write18{find "#1" -name "*.tex" 2>/dev/null | sort -V | sed 's/.*/\\input{&}/' > "\apuntesentradasfile"}%
    \input{\apuntesentradasfile}%
  \endgroup
}
